\section{The order-three Hodge representation}

The number eleven in the cyclic cubic target is forced by the Hodge
representation; it is not merely the difference between thirty and the base
rank eight.

Let $S_C$ be the cubic genus-one pullback.  It has
\[
\chi(\mathcal O_{S_C})=3,
\qquad
q=1,
\qquad
p_g=3,
\qquad
e=36.
\]
Hence
\[
\boxed{b_2=38,\qquad h^{1,1}=32.}
\]

Let $\sigma$ generate the deck group.  Its fixed locus is the union of the
three smooth elliptic fibres above the branch points, whose Euler number is
zero.  On $H^1$ and $H^3$, the trace is $-1$, because the quotient of the
base curve is $\mathbf P^1$.  The topological Lefschetz formula gives
\[
\operatorname{Tr}(\sigma|H^2)=-4.
\]
Writing the dimensions of the three complex character spaces as
\[
n_0,n,n,
\]
one has
\[
n_0+2n=38,
\qquad
n_0-n=-4.
\]
Therefore
\[
\boxed{(n_0,n,n)=(10,14,14).}
\]

The holomorphic two-forms have character multiplicities
\[
(0,2,1)
\]
up to exchanging the two nontrivial characters.  Complex conjugation reverses
the last two entries for $H^{0,2}$.  It follows that
\[
\boxed{
\dim H^{1,1}_1=10,
\qquad
\dim H^{1,1}_{\zeta_3}
=
\dim H^{1,1}_{\zeta_3^2}=11.
}
\]

Since all singular fibres remain irreducible,
\[
\rank\MW(S_C)=\rho(S_C)-2.
\]
The invariant Mordell--Weil lattice is $E_8(3)$, and each nontrivial character
has capacity eleven.  Thus the following are equivalent:
\[
\boxed{
\begin{aligned}
&\rank\MW(S_C)=30,\\
&\rho(S_C)=32=h^{1,1}(S_C),\\
&\dim\MW_{\zeta_3}=\dim\MW_{\zeta_3^2}=11,\\
&\text{all nontrivial-character }(1,1)\text{ classes are algebraic.}
\end{aligned}
}
\]

A rank-thirty cubic pullback is therefore a Picard-maximal irregular elliptic
surface with a trace-zero Eisenstein Mordell--Weil lattice of complex rank
eleven.  This connects the explicit trisection search to order-three
Noether--Lefschetz loci, CM points, Delsarte surfaces, and ternary lattice
gluing.
